\scanpage{48}
\mspage{4)}
% archive sheet number: 47
On a un foncteur
\[
\Gamma:\quad S_*\mathscr{T}\unsure{\mathbb D|D}_{\mathrm{top}}^{\pm}\longrightarrow\deleted{\mathrm{MD}}(\mathrm{Grabipond})\tag*{(*)}
\]
qui associe à un objet \inserted{$(X,\widetilde\Sigma,R)$} de $S_*\mathscr{T}\mathbb D_{\mathrm{top}}^{\pm}$ le graphe \inserted{MD} $(S,\vec A,\sigma,\deleted{\unsure{r|\rho}}\inserted{o},\deleted{\unsure{g}}\unsure{\gamma|g},\rho)$ suivant:
\[
S=\pi_0(X\deleted{\unsure{,}}\setminus\widetilde\Sigma)\simeq\pi_0(\mathrm{Déc}(X,\widetilde\Sigma))
\]
$\vec A={}$ ens des couples $(C,\deleted{\unsure{g}}\omega)$, où $C\deleted{\in A}\inserted{\in\pi_0}(\widetilde\Sigma)$, et où $\omega\in\mathrm{Or}(C)\inserted{\simeq\mathrm{rive}_X(C)}$, soumis à la condition que si $C\in\pi_0(\partial X)$, alors $\omega$ est l'orientation induite i.e. la rive qui se trouve dans $X$

$\sigma={}$ échange des rives (quand on a le droit d'échanger, sinon, i.e. si $C\in\pi_0(\partial X)$, on pose $\sigma(a)=a$)

$\deleted{\unsure{R|\rho}}\inserted{o}:{}$ associe à $(C,\text{rive de }C)$ l'unique ``\unsure{face}'' $\in\pi_0(X\setminus\widetilde\Sigma)$ qui contient cette rive

\[
\gamma:\underset{\Vert\atop\pi_0(X\setminus\widetilde\Sigma)}{S}\longrightarrow\mathbb N\qquad\text{application \uline{genre} des composants }X_i.
\]
\[
\unsure{r|\rho}:\quad A=\pi_0(\widetilde\Sigma)\longrightarrow\mathbb N:\quad C\longmapsto\mathrm{Card}(R\cap C).
\]

Le foncteur (*) \deleted{est \unsure{ess.} surjectif et}\inserted{induit un \uline{iso} sur les} $\pi_0$ et un surjectif sur les $\pi_1$ (c'est un \uline{$0$-isomorphisme}\deleted{\unsure{,}}, en langage homotopique).

\deleted{\unsure{Prop.}}:

\uline{\unsure{Remarque|Remarques}} Dans $S_*\mathscr{T}\mathbb D_{\mathrm{top}}^{\pm}$ il y a une opération évidente de \uline{somme} \inserted{\uline{discrète}}, dans \deleted{\mathrm{MD}}\inserted{$\mathrm{Grabipond}$} aussi, et (*) y commute.
